\chapter*{\centering ABSTRACT}
\addcontentsline{toc}{chapter}{ABSTRACT}

\begin{center}
\setstretch{1.2}
\begin{tabular}{l l}
Name          & : Gilang Pratama Putra Siswanto \\ 
Study Program & : Physics \\ 
Title         & : \begin{tabular}[c]{@{}l@{}}Design of an Automation Experiment Kit \\for Wien’s Displacement Law Based on TCS34725 \\ RGB Sensor, Arduino Uno R3, and Python 3 Interface\end{tabular} \\ 
\end{tabular}
\end{center}

\vspace{1cm}

\begin{spacing}{1.25}
\noindent This research aims to analyze the accuracy and precision of the Wien Displacement Law experimental kit based on Arduino Uno and the TCS34725 sensor in measuring the relationship between temperature and the wavelength of black body radiation. By adopting modern technology, this experiment was conducted at the Instrumentation and Robotics Laboratory of the State Islamic University of Sunan Gunung Djati from September to October 2024. The research methodology involves the use of temperature sensors and spectrum sensors integrated with a microcontroller to determine the maximum wavelength emitted by an object at a certain temperature. The testing results indicate that the developed digital system operates optimally and accurately, producing a Wien displacement constant (b) value of 0.002897773 with a precision of 1,8796E-09 and an accuracy of 99.99\% compared to the literature value. This research is expected to provide a deeper understanding of the application of automation in physics experiments and its contribution to the development of science. \\

\noindent\textbf{Keywords:} Wien Displacement Law, Arduino Uno, TCS34725 sensor, temperature measurement, black body radiation.
\end{spacing}
